\section{四因素压力测试设置}
\label{sec:03_06}

初始合成条件在固定资源数下改变任务数或任务阶数，会同时改变资源竞争程度。为进一步区分这些作用，本节建立四因素压力测试：分别控制候选任务数、单项任务的资源需求数、同一资源关联的任务数以及收益跨度，再在公开拓扑上实施对应的结构或权重干预。前一组用于解释可控生成机制下的变化，后一组用于观察非规则关联结构中的表现。两组均沿用单位容量的任务选择模型，但公开拓扑的干预不一定只改变一个实际结构指标，因而需要同时记录干预名称与生成后的实际参数。

\subsection{四个因素及其关联规模}

设资源集合为$V$，候选任务索引集合为$E$，任务$e\in E$需要的非空资源集合为$R_e\subseteq V$。关联矩阵$A\in\{0,1\}^{m\times n}$的元素为$A_{ve}=\mathbb{I}_{\{v\in R_e\}}$，其中$n=|E|$为候选任务数，$m=|V|$为资源数。每项任务的选取变量$x_e\in\{0,1\}$表示是否执行该任务，固定正整数$w_e$表示其无量纲收益。全部压力条件求解同一形式的问题：
\begin{equation}
 \begin{aligned}
 \max_x\quad &\sum_{e\in E}w_ex_e,\\
 \text{s.t.}\quad &\sum_{e\in E}A_{ve}x_e\leq1,
 &&\forall v\in V,\\
 &x_e\in\{0,1\},&&\forall e\in E.
 \end{aligned}
 \label{s0306:model}
\end{equation}
容量右端的$1$表示一个不可分资源单位；这些收益和资源计数均不对应实际飞行能量或货币。

四个控制因素记为$n,r,\Delta,W$。其中$n$控制二元变量数，$r$控制任务资源需求规模，$\Delta$控制资源竞争规模，$W$控制生成权重的跨度。任务阶数、资源度及总关联项数分别定义为
\begin{align}
 r_e&=|R_e|,&&e\in E,\label{s0306:edge_rank}\\
 d_v&=\sum_{e\in E}A_{ve},&&v\in V,\label{s0306:resource_degree}\\
 r&=\max_{e\in E}r_e,\qquad
 \Delta=\max_{v\in V}d_v,\qquad
 \operatorname{nnz}(A)=\sum_{e\in E}r_e=\sum_{v\in V}d_v.
 \label{s0306:incidence_count}
\end{align}
$r_e,d_v,r,\Delta$及$\operatorname{nnz}(A)$均为无量纲正整数。这里没有神经特征维数这一因素，不能把任务阶数或资源度解释成学习模型的特征宽度。

在规则生成器中，每项任务恰好需要$r$种资源，每种资源恰好关联$\Delta$项任务。因此，关联矩阵的行列计数满足
\begin{equation}
 \operatorname{nnz}(A)=nr=m\Delta,
 \qquad m=\frac{nr}{\Delta}.
 \label{s0306:regular_counts}
\end{equation}
这说明$m$不是能够同时固定的第五个独立控制量。增加$n$或$r$会增加资源数和关联项数；增加$\Delta$则在保持$n,r$时减少资源数，同时把更多任务放入同一容量约束。解释计算开销时，必须同时查看$n,m,r,\Delta$和$\operatorname{nnz}(A)$，不能只报告被改变的一个因素。

\subsection{规则超图的分层构造与配对关系}

规则条件采用以下水平集合：
\begin{equation}
 \begin{aligned}
 n&\in\{256,1024,4096,16384\},\\
 r&\in\{4,8,16,32\},\\
 \Delta&\in\{4,8,16,32\},\\
 W&\in\{1,10,100,1000\}.
 \end{aligned}
 \label{s0306:levels}
\end{equation}
基准配置为$(n,r,\Delta,W)=(256,4,4,100)$。每次只把一个控制因素改为其余三个水平，另外三个控制因素保持基准值。因此配置数为$1+4\times3=13$，而不是四个水平集合的笛卡尔积。每个配置使用结构种子$s\in\{0,1,2\}$，共形成$13\times3=39$个规则条件。该设计可以观察基准附近的单因素变化，但没有完整覆盖四因素之间的交互作用。

下面给出生成器的组合含义。将任务编号固定为$E_n=\{0,\ldots,n-1\}$。对每个层索引$\ell\in\{0,\ldots,r-1\}$，生成$E_n$的一个随机排列$\pi_{s,\ell,n}$。所有取值均满足$\Delta\mid n$，故可把排列按连续$\Delta$个任务划分为$n/\Delta$组。第$\ell$层、第$k$组的任务集合为
\begin{equation}
 P_{\ell,k}^{(s)}=
 \{\pi_{s,\ell,n}(k\Delta+t):t=0,\ldots,\Delta-1\},
 \quad
 \ell=0,\ldots,r-1,\quad k=0,\ldots,n/\Delta-1.
 \label{s0306:partition}
\end{equation}
每组建立一个容量为$1$的资源$v_{\ell,k}$，组内任务共同占用该资源。于是
\begin{align}
 V&=\{v_{\ell,k}:0\leq\ell<r,\ 0\leq k<n/\Delta\},
 \label{s0306:layer_resources}\\
 R_j&=\{v_{\ell,k}:j\in P_{\ell,k}^{(s)}\},
 &&j\in E_n.\label{s0306:layer_incidence}
\end{align}
同层分组构成任务集合的一个划分，所以每项任务在每层恰好关联一个资源；不同层的资源身份不同。因此$|R_j|=r$，且每个$v_{\ell,k}$恰好关联$\Delta$项任务，式\eqref{s0306:regular_counts}随之成立。相应容量约束可以逐层写为
\begin{equation}
 \sum_{j\in P_{\ell,k}^{(s)}}x_j\leq1,
 \qquad
 \ell=0,\ldots,r-1,\quad k=0,\ldots,n/\Delta-1.
 \label{s0306:partition_constraints}
\end{equation}
它限制同一资源组内至多接受一项任务；同一任务在所有层中的约束必须同时满足。

收益使用与层分区分离的伪随机流。令$\rho_{s,n}$为$\{0,\ldots,n-1\}$的一个随机排列，定义
\begin{equation}
 q_j=\frac{\rho_{s,n}(j)}{n-1},\qquad
 w_j=\max\{1,\operatorname{round}(W^{q_j})\},
 \qquad j\in E_n.
 \label{s0306:regular_weights}
\end{equation}
$q_j\in[0,1]$为无量纲分位位置，$W\geq1$为无量纲跨度参数，$\operatorname{round}$按实现采用的最近整数舍入规则取整。分位位置集合恰为$\{0,1/(n-1),\ldots,1\}$；随机性只改变这些位置分配给哪项任务。对本节的整数$W$，实际最小权重为$1$、最大权重为$W$，故$w_{\max}/w_{\min}=W$。当$W=1$时全部任务等权，目标退化为最大化接受任务数。

实现通过NumPy的\texttt{SeedSequence}构造分离的随机流：第$\ell$层使用种子输入$[s,\ell,7741]$，收益排列使用$[s,9127]$。这里的固定整数用于区分随机流，不是模型参数。结构随机流与收益随机流分离，使固定$n,s$时改变$r$或$\Delta$不会推进收益排列，改变$W$也不会改变资源分区。

配对关系须按生成机制分别理解。在固定$n,\Delta,s$时增加$r$，已有层保持不变，只增加新的资源层，故原任务的资源需求按层嵌套。在固定$n,r,s$时，$\Delta$从$4$增至$8$、$16$或$32$使用相同层排列，把相邻小组依次合并成大组。例如，容量为$1$的两个小组被一个容量为$1$的大组替代后，原先可以分别在两组各选一项任务，现在只能在合并组内总共选一项。令$\mathcal F_{n,r,\Delta,s}\subseteq\{0,1\}^n$表示规则条件的可行域，则在其余索引固定且水平满足上述嵌套关系时，有
\begin{align}
 \mathcal F_{n,r_2,\Delta,s}&\subseteq
 \mathcal F_{n,r_1,\Delta,s},&&r_2\geq r_1,
 \label{s0306:rank_nesting}\\
 \mathcal F_{n,r,\Delta_2,s}&\subseteq
 \mathcal F_{n,r,\Delta_1,s},&&\Delta_2\geq\Delta_1.
 \label{s0306:degree_nesting}
\end{align}
第二个包含关系依赖这里按同一排列进行整倍数粗化的构造，不能推广成任意两个最大资源度不同的超图之间的包含关系。两式也不意味着启发式输出或求解时间必须单调变化。

改变$n$则没有同样的前缀保证。虽然使用相同的种子输入，长度不同的随机排列会改变任务分区，收益排列及分位位置也会改变。因此，$n=256$的任务资源关系不能声称是$n=1024$条件的原样前缀。这一点不同于第\ref{sec:03_05}节的初始合成构造。同一个$s$在本节用于组织配对条件，但并不保证跨所有因素、所有规模的每条任务资源关联都相同；$39$个条件也不能当作$39$个相互独立的随机结构样本。

\subsection{公开拓扑的嵌套抽取与干预含义}

另一组条件从第\ref{sec:03_05}节的完整\texttt{superblue19}拓扑派生。该拓扑分别用于JGAA公开输入与ESA代码权重重建条件\cite{GrossmannJRS2026,ReinstadtlerFPU2025}。这里的派生输入经过抽取或修改，应与保持完整拓扑的公开基准区分。

抽取以关联任务最多的资源为起点；若存在并列，按资源标识的字典序选择。先将该资源关联的任务加入队列，再沿“任务---资源---任务”的关联关系作广度优先遍历，任务与资源均按实现规定避免重复访问。设得到的任务发现序列为$(e_1,\ldots,e_{65536})$，取
\begin{equation}
 E_N^{\mathrm{sub}}=\{e_1,\ldots,e_N\},
 \qquad N\in\{4096,16384,65536\}.
 \label{s0306:public_subsets}
\end{equation}
于是三个任务集合具有包含关系。写入求解输入时，集合内任务恢复为源文件中的任务顺序，而非按广度优先发现顺序输入；保留任务的全部原有关联，并删除不再被任何任务使用的资源。这样的抽取围绕一个高竞争资源展开，具有明确的选择偏向，不是对完整拓扑的均匀随机抽样。

每个抽取集合采用两种收益。第一种称为\texttt{random100}，继承源档案的随机整数任务权重；第二种称为LHG，使用完整源拓扑的最大任务阶数$r_{\mathrm{full}}$，令
\begin{equation}
 w_e^{\mathrm{LHG}}=r_{\mathrm{full}}+1-|R_e^{\mathrm{full}}|,
 \qquad e\in E_N^{\mathrm{sub}}.
 \label{s0306:public_lhg}
\end{equation}
其中$R_e^{\mathrm{full}}$为任务在完整源拓扑中的资源集合。$r_{\mathrm{full}}$不随子集规模重新计算；同一任务在三个嵌套子集中保持相同收益。在后续截断或拆分关联时，该收益也不重新按修改后的阶数计算。因此，LHG表示干预前的赋权来源，不表示每个派生输入都重新满足“当前最大阶数加一减当前任务阶数”的关系。

对每种收益，设置三个任务规模条件，并在$N=16384$的子集上设置两种阶数截断、两种资源度截断、两种收益跨度变换和一种整体收益缩放，共$3+2+2+2+1=10$个条件。两种收益合计$20$个条件，其中$18$个为结构或收益跨度干预，另有$2$个整体尺度控制。其操作及数学含义如下。

\paragraph{截断任务阶数。}
对每项任务，依据任务标识与资源标识组成的字符串计算SHA-256摘要，以摘要排序给出确定的资源优先序，再保留前$c$个关联，其中$c\in\{8,32\}$。优先序不依赖求解结果。若$R_e^{(c)}$表示截断后的资源集合，则
\begin{equation}
 R_e^{(c)}\subseteq R_e^{\mathrm{sub}},\qquad
 |R_e^{(c)}|=\min\{|R_e^{\mathrm{sub}}|,c\},
 \qquad e\in E_{16384}^{\mathrm{sub}}.
 \label{s0306:rank_cap}
\end{equation}
该操作保持任务数与收益，但删除了部分资源占用要求。因而原可行选择仍然可行，新的可行域可能扩大。资源的关联任务数也可能下降，未被使用的资源会被移除，故它不是只改变$r$而严格保持$m$、$\Delta$及关联总数的干预。资源优先序固定还使$c=8$的保留集合包含于$c=32$的保留集合。

\paragraph{拆分高关联度资源。}
对每个原资源$v$，按输入任务顺序列出使用该资源的任务，并将其连续划分为至多$c$项任务一组，其中$c\in\{32,128\}$。每组建立一个独立的单位容量资源。记原任务集合为$E(v)$，划分后的集合为$E_{v,1},\ldots,E_{v,k_v}$，其中$k_v=\lceil|E(v)|/c\rceil$。原约束及替代约束分别为
\begin{align}
 \sum_{e\in E(v)}x_e&\leq1,
 \label{s0306:before_split}\\
 \sum_{e\in E_{v,k}}x_e&\leq1,
 &&k=1,\ldots,k_v.
 \label{s0306:after_split}
\end{align}
划分集合两两不交，且并集为$E(v)$。满足式\eqref{s0306:before_split}必然满足式\eqref{s0306:after_split}，反向关系一般不成立：来自不同组的任务现在可以分别占用不同资源。因此，拆分资源放松了原来的冲突限制，改变了优化问题的可行域。每个原关联被替换为一个新关联，故$n$、各任务阶数和$\operatorname{nnz}(A)$保持不变，资源数则可能增加，新资源度不超过$c$。

\paragraph{改变收益跨度。}
在$16384$任务子集中，先把不同的原权重排序为$a_0<\cdots<a_{K-1}$，其中$K$为不同权重的个数。对原权重为$a_i$的任务定义
\begin{equation}
 \eta_e=\frac{i}{\max\{1,K-1\}},\qquad
 w'_e=\max\{1,\operatorname{round}(W^{\eta_e})\},
 \quad W\in\{1,1000\}.
 \label{s0306:public_spread}
\end{equation}
这里$\eta_e\in[0,1]$按不同权重值的顺序确定，不按每个权重出现的频数确定；原权重相同的任务仍得到相同新权重。该变换保持关联与可行域，但改变收益间的相对差异，整数舍入还可能合并部分权重等级。特别是$W=1$使全部任务等权，所以它不能解释为只改变目标函数的计量单位。

\paragraph{整体收益缩放。}
尺度控制保留所有关联，令$w'_e=10w_e$。若$\mathcal F$为固定的二元可行域，则
\begin{equation}
 \underset{x\in\mathcal F}{\operatorname{arg\,max}}
 \sum_e10w_ex_e
 =
 \underset{x\in\mathcal F}{\operatorname{arg\,max}}
 \sum_ew_ex_e.
 \label{s0306:scale_invariance}
\end{equation}
正比例缩放使所有可行解的收益同时乘以$10$，保持整数最优选择集合。这一数学性质不保证具体实现的访问顺序、数值容差处理、停止轨迹和运行时间完全相同，故尺度控制仍需独立记录运行结果。

这些公开派生条件共享同一来源拓扑，两个收益体系共享关联结构，三个规模子集又相互嵌套。映射表中的种子$0$是该组条件的记录字段，不能解释为完成了多个独立拓扑随机重复。因此，$20$是条件数，不是独立数据集数。

\subsection{全部条件的标识与实际参数}

表\ref{tab:stress_condition_map}给出全部$59$个条件的编号、来源、干预轴、水平和种子。表\ref{tab:stress_parameters}使用同一S编号逐行给出生成后的实际结构参数，两个表通过编号一一对应。输入名称中的后缀只用于标识条件，不另行定义数学变量。对于基准行，水平栏中的横线表示使用完整基准配置，并不表示参数缺失。

{\small\setlength{\tabcolsep}{5pt}
\begin{longtable}{lllrl}
\caption{条件编号、来源与干预的对应}\label{tab:stress_condition_map}\\
\toprule 编号 & 来源 & 干预轴 & 水平 & 种子\\\midrule
\endfirsthead
\toprule 编号 & 来源 & 干预轴 & 水平 & 种子\\\midrule
\endhead
\midrule\multicolumn{5}{r}{续下页}\\\endfoot
\bottomrule\endlastfoot
S01 & 规则超图 & 基准 & -- & 0\\
S02 & 规则超图 & 资源度 & 16 & 0\\
S03 & 规则超图 & 资源度 & 32 & 0\\
S04 & 规则超图 & 资源度 & 8 & 0\\
S05 & 规则超图 & 任务阶数 & 16 & 0\\
S06 & 规则超图 & 任务阶数 & 32 & 0\\
S07 & 规则超图 & 任务阶数 & 8 & 0\\
S08 & 规则超图 & 收益跨度 & 1 & 0\\
S09 & 规则超图 & 收益跨度 & 10 & 0\\
S10 & 规则超图 & 收益跨度 & 1000 & 0\\
S11 & 规则超图 & 基准 & -- & 1\\
S12 & 规则超图 & 资源度 & 16 & 1\\
S13 & 规则超图 & 资源度 & 32 & 1\\
S14 & 规则超图 & 资源度 & 8 & 1\\
S15 & 规则超图 & 任务阶数 & 16 & 1\\
S16 & 规则超图 & 任务阶数 & 32 & 1\\
S17 & 规则超图 & 任务阶数 & 8 & 1\\
S18 & 规则超图 & 收益跨度 & 1 & 1\\
S19 & 规则超图 & 收益跨度 & 10 & 1\\
S20 & 规则超图 & 收益跨度 & 1000 & 1\\
S21 & 规则超图 & 基准 & -- & 2\\
S22 & 规则超图 & 资源度 & 16 & 2\\
S23 & 规则超图 & 资源度 & 32 & 2\\
S24 & 规则超图 & 资源度 & 8 & 2\\
S25 & 规则超图 & 任务阶数 & 16 & 2\\
S26 & 规则超图 & 任务阶数 & 32 & 2\\
S27 & 规则超图 & 任务阶数 & 8 & 2\\
S28 & 规则超图 & 收益跨度 & 1 & 2\\
S29 & 规则超图 & 收益跨度 & 10 & 2\\
S30 & 规则超图 & 收益跨度 & 1000 & 2\\
S31 & 规则超图 & 任务数 & 1024 & 0\\
S32 & 规则超图 & 任务数 & 1024 & 1\\
S33 & 规则超图 & 任务数 & 1024 & 2\\
S34 & 规则超图 & 任务数 & 4096 & 0\\
S35 & 公开拓扑/阶数重加权 & 任务数 & 4096 & 0\\
S36 & 公开拓扑/随机权重 & 任务数 & 4096 & 0\\
S37 & 规则超图 & 任务数 & 4096 & 1\\
S38 & 规则超图 & 任务数 & 4096 & 2\\
S39 & 规则超图 & 任务数 & 16384 & 0\\
S40 & 公开拓扑/阶数重加权 & 拆分资源 & 128 & 0\\
S41 & 公开拓扑/阶数重加权 & 拆分资源 & 32 & 0\\
S42 & 公开拓扑/阶数重加权 & 任务数 & 16384 & 0\\
S43 & 公开拓扑/阶数重加权 & 截断阶数 & 32 & 0\\
S44 & 公开拓扑/阶数重加权 & 截断阶数 & 8 & 0\\
S45 & 公开拓扑/阶数重加权 & 整体权重倍数 & 10 & 0\\
S46 & 公开拓扑/阶数重加权 & 收益跨度 & 1 & 0\\
S47 & 公开拓扑/阶数重加权 & 收益跨度 & 1000 & 0\\
S48 & 公开拓扑/随机权重 & 拆分资源 & 128 & 0\\
S49 & 公开拓扑/随机权重 & 拆分资源 & 32 & 0\\
S50 & 公开拓扑/随机权重 & 任务数 & 16384 & 0\\
S51 & 公开拓扑/随机权重 & 截断阶数 & 32 & 0\\
S52 & 公开拓扑/随机权重 & 截断阶数 & 8 & 0\\
S53 & 公开拓扑/随机权重 & 整体权重倍数 & 10 & 0\\
S54 & 公开拓扑/随机权重 & 收益跨度 & 1 & 0\\
S55 & 公开拓扑/随机权重 & 收益跨度 & 1000 & 0\\
S56 & 规则超图 & 任务数 & 16384 & 1\\
S57 & 规则超图 & 任务数 & 16384 & 2\\
S58 & 公开拓扑/阶数重加权 & 任务数 & 65536 & 0\\
S59 & 公开拓扑/随机权重 & 任务数 & 65536 & 0\\
\end{longtable}}


{\small\setlength{\tabcolsep}{4pt}
\begin{longtable}{lrrrrrr}
\caption{压力测试全部条件的实际结构参数}\label{tab:stress_parameters}\\
\toprule
编号 & $n$ & $m$ & $r$ & $\Delta$ & $\operatorname{nnz}(A)$ & $w_{\max}/w_{\min}$\\
\midrule
\endfirsthead
\toprule
编号 & $n$ & $m$ & $r$ & $\Delta$ & $\operatorname{nnz}(A)$ & $w_{\max}/w_{\min}$\\
\midrule
\endhead
\midrule
\multicolumn{7}{r}{续下页}\\
\endfoot
\bottomrule
\endlastfoot
S01 & 256 & 256 & 4 & 4 & 1024 & 100\\
S02 & 256 & 64 & 4 & 16 & 1024 & 100\\
S03 & 256 & 32 & 4 & 32 & 1024 & 100\\
S04 & 256 & 128 & 4 & 8 & 1024 & 100\\
S05 & 256 & 1024 & 16 & 4 & 4096 & 100\\
S06 & 256 & 2048 & 32 & 4 & 8192 & 100\\
S07 & 256 & 512 & 8 & 4 & 2048 & 100\\
S08 & 256 & 256 & 4 & 4 & 1024 & 1\\
S09 & 256 & 256 & 4 & 4 & 1024 & 10\\
S10 & 256 & 256 & 4 & 4 & 1024 & 1000\\
S11 & 256 & 256 & 4 & 4 & 1024 & 100\\
S12 & 256 & 64 & 4 & 16 & 1024 & 100\\
S13 & 256 & 32 & 4 & 32 & 1024 & 100\\
S14 & 256 & 128 & 4 & 8 & 1024 & 100\\
S15 & 256 & 1024 & 16 & 4 & 4096 & 100\\
S16 & 256 & 2048 & 32 & 4 & 8192 & 100\\
S17 & 256 & 512 & 8 & 4 & 2048 & 100\\
S18 & 256 & 256 & 4 & 4 & 1024 & 1\\
S19 & 256 & 256 & 4 & 4 & 1024 & 10\\
S20 & 256 & 256 & 4 & 4 & 1024 & 1000\\
S21 & 256 & 256 & 4 & 4 & 1024 & 100\\
S22 & 256 & 64 & 4 & 16 & 1024 & 100\\
S23 & 256 & 32 & 4 & 32 & 1024 & 100\\
S24 & 256 & 128 & 4 & 8 & 1024 & 100\\
S25 & 256 & 1024 & 16 & 4 & 4096 & 100\\
S26 & 256 & 2048 & 32 & 4 & 8192 & 100\\
S27 & 256 & 512 & 8 & 4 & 2048 & 100\\
S28 & 256 & 256 & 4 & 4 & 1024 & 1\\
S29 & 256 & 256 & 4 & 4 & 1024 & 10\\
S30 & 256 & 256 & 4 & 4 & 1024 & 1000\\
S31 & 1024 & 1024 & 4 & 4 & 4096 & 100\\
S32 & 1024 & 1024 & 4 & 4 & 4096 & 100\\
S33 & 1024 & 1024 & 4 & 4 & 4096 & 100\\
S34 & 4096 & 4096 & 4 & 4 & 16384 & 100\\
S35 & 4096 & 5808 & 269 & 1507 & 10962 & 1.03\\
S36 & 4096 & 5808 & 269 & 1507 & 10962 & 100\\
S37 & 4096 & 4096 & 4 & 4 & 16384 & 100\\
S38 & 4096 & 4096 & 4 & 4 & 16384 & 100\\
S39 & 16384 & 16384 & 4 & 4 & 65536 & 100\\
S40 & 16384 & 33966 & 8618 & 128 & 62561 & 17.015\\
S41 & 16384 & 34038 & 8618 & 32 & 62561 & 17.015\\
S42 & 16384 & 33944 & 8618 & 1507 & 62561 & 17.015\\
S43 & 16384 & 22352 & 32 & 1507 & 46876 & 17.015\\
S44 & 16384 & 22062 & 8 & 1507 & 44684 & 17.015\\
S45 & 16384 & 33944 & 8618 & 1507 & 62561 & 17.015\\
S46 & 16384 & 33944 & 8618 & 1507 & 62561 & 1\\
S47 & 16384 & 33944 & 8618 & 1507 & 62561 & 1000\\
S48 & 16384 & 33966 & 8618 & 128 & 62561 & 100\\
S49 & 16384 & 34038 & 8618 & 32 & 62561 & 100\\
S50 & 16384 & 33944 & 8618 & 1507 & 62561 & 100\\
S51 & 16384 & 22352 & 32 & 1507 & 46876 & 100\\
S52 & 16384 & 22062 & 8 & 1507 & 44684 & 100\\
S53 & 16384 & 33944 & 8618 & 1507 & 62561 & 100\\
S54 & 16384 & 33944 & 8618 & 1507 & 62561 & 1\\
S55 & 16384 & 33944 & 8618 & 1507 & 62561 & 1000\\
S56 & 16384 & 16384 & 4 & 4 & 65536 & 100\\
S57 & 16384 & 16384 & 4 & 4 & 65536 & 100\\
S58 & 65536 & 105235 & 8618 & 1507 & 235201 & 17.017\\
S59 & 65536 & 105235 & 8618 & 1507 & 235201 & 100\\
\end{longtable}
}



两表需要合并阅读。例如，S43是LHG赋权的公开子集在阶数上限$32$下的干预，其实际$n=16384$、$m=22352$、$r=32$、$\Delta=1507$及$\operatorname{nnz}(A)=46876$来自同一个输入。S41则采用资源拆分上限$32$，实际$r=8618$、$\Delta=32$及$\operatorname{nnz}(A)=62561$。前者删除关联，后者替换资源身份；即使干预水平都写为$32$，也不代表同一种结构变化。参数表中的收益比反映实际输入，公开LHG子集的收益比不必等于规则条件的$W$。

输入核验检查任务标识和资源标识的唯一性、资源容量均为$1$、任务资源集合非空且无重复成员、所有资源引用存在，以及权重为正整数。还检查全部任务权重之和小于$2147483647$，以满足原作者实现使用的整数目标范围。规则条件另外核对精确阶数、精确资源度和$\operatorname{nnz}(A)=nr$。这些检查保证各方法接收的模型对象一致；它们本身不构成求解质量结论。

\subsection{求解方法、停止条件与计时边界}

每个条件均运行Streaming、Greedy、ILS1000和SCIP，合计$59\times4=236$次。原始结果中的任务选择已按对应原输入重新计算资源占用与收益，$236$次输出均通过可行性核验。核验不通过修改任务集合来修复输出；SCIP另外核对重新计算的收益与求解器报告收益的一致性。该覆盖范围是压力测试自身的运行数，不包括初始实验或原作者完整流程的运行。

SCIP采用10.0.2版本及PySCIPOpt~6.2.1接口，按式\eqref{s0306:model}建立整数规划。每次仅提供全零可行解作为起点，不使用Streaming、Greedy、ILS或历史SCIP选择作热启动。随机种子偏移设为$0$，线程参数设为$1$，相对间隙停止目标为$5\%$，求解器内存参数为$1536\,\mathrm{MiB}$，不另设求解时限。无求解时限并不意味着必然证明最优：运行仍可能因达到间隙目标、内存限制或其他求解状态而结束。因此必须保留实际停止状态、最终可行值和上界，不能把停止参数当作每次运行均达标的证明。

Streaming采用$\delta_g$配置，参数$\varepsilon=0.1$，完成一次给定任务顺序的扫描及原有输出过程后结束\cite{ReinstadtlerFPU2025}。Greedy完成贪心构造后结束。ILS1000在贪心构造基础上执行原作者实现的局部搜索，其$1000$表示连续未改进计数的停止阈值：发现改进后计数重置，代码在该计数大于$1000$时停止\cite{GrossmannJRS2026}。它不是固定执行$1000$次总迭代。压力测试中的Greedy和ILS1000调用固定算法种子$0$；规则输入的三个结构种子用于改变输入，不是同一输入上的三个算法随机重复。这里也没有把归约、展开及ILS15组成的原作者完整流程混入四方法比较。

设条件$c$上方法$a$的总计时时长为$T_{c,a}>0$，单位为$\mathrm{s}$，其口径为
\begin{equation}
 T_{c,a}=T_{c,a}^{\mathrm{build}}
       +T_{c,a}^{\mathrm{algorithm}}
       +T_{c,a}^{\mathrm{output}}.
 \label{s0306:timing}
\end{equation}
三个分量分别表示建模或输入转换、算法执行、结果提取；它们均为相应计时区间的墙钟时间。磁盘输入读取、外部可行性核验和结果文件写入不并入这一总段，不能把结果提取时间理解为完整磁盘输出成本。进程墙钟时间和CPU时间另行记录，前者还覆盖启动及外围工作，后者衡量进程消耗的处理器时间，均不与式\eqref{s0306:timing}混用。

执行队列最多同时安排两个SCIP进程和一个源码方法进程，源码方法按Streaming、Greedy、ILS1000的次序分别遍历条件；已有无关求解进程保持运行。因此，单次SCIP的单线程设置不等于整台计算机仅有一个活动计算任务，墙钟时间可能受到并发争用影响。监控每隔$1\,\mathrm{s}$采样一次进程驻留内存，并设有$2560\,\mathrm{MiB}$的进程内存保护阈值。采样得到的是观测到的驻留内存峰值，不能当作精确的瞬时分配峰值，也不能与SCIP内部的$1536\,\mathrm{MiB}$内存参数直接等同。

四种方法的停止机制不同：SCIP承担搜索和界收敛，Streaming完成单遍过程，Greedy完成构造，ILS按连续未改进次数结束。因而这些时间可以描述各自协议下取得相应输出的成本，但不构成等计算预算的算法速度排名，也不能与初始$10\,\mathrm{s}$实验或原作者完整流程的计时直接合并。

\subsection{质量参照、联合判据与统计单位}

对条件$c$，记SCIP最终输出经重新核验的可行收益为$L_c$，其他方法$a$的可行收益为$h_{c,a}$，该条件未知的最优收益为$z_c^*$。若SCIP还提供有限有效上界$U_c$，则
\begin{equation}
 0\leq L_c\leq z_c^*\leq U_c,
 \qquad 0\leq h_{c,a}\leq z_c^*.
 \label{s0306:bounds}
\end{equation}
这些量均采用该条件的无量纲收益尺度。$L_c$是已取得的可行值，不因被用作比较基准就成为最优值；$h_{c,a}>L_c$也可能发生，表示该方法找到了比本次SCIP输出收益更高的可行选择。

当$L_c>0$且计时有效时，定义相对可行值比和时间比为
\begin{equation}
 Q_{c,a}=\frac{h_{c,a}}{L_c},\qquad
 R_{c,a}=\frac{T_{c,\mathrm{SCIP}}}{T_{c,a}}.
 \label{s0306:comparison_ratios}
\end{equation}
二者均无量纲。预定联合判据要求同一个条件、同一次方法输出同时满足
\begin{equation}
 Q_{c,a}\geq0.95,
 \qquad R_{c,a}\geq2.
 \label{s0306:joint_criterion}
\end{equation}
第一项要求收益至少达到SCIP已得可行值的$95\%$；第二项要求在所定义的总计时口径下，SCIP时间至少为该方法的两倍。不能把某个条件的质量比与另一条件的时间比拼接后判定支持，也不能只用跨种子均值判断所有种子是否同时满足要求。

对规则生成器的某个固定配置$g$，令$c(g,s)$表示结构种子$s\in\{0,1,2\}$对应的条件。方法$a$的联合支持数定义为
\begin{equation}
 K_{g,a}=\sum_{s\in\{0,1,2\}}
 \mathbb{I}_{\{Q_{c(g,s),a}\geq0.95\}}
 \mathbb{I}_{\{R_{c(g,s),a}\geq2\}},
 \qquad K_{g,a}\in\{0,1,2,3\}.
 \label{s0306:support_count}
\end{equation}
每个配置均报告支持数相对于预定三个种子的覆盖，同时保留各比值的均值、最小值和最大值。缺少有效输出、核验失败或比值分母无效时，不能将该条件记为支持；应保留原因及计划覆盖数。实际低于质量阈值、未达到时间阈值或者两项均不满足的结果，也不从配置中剔除。公开派生条件没有三个独立结构种子的对应重复，故不能套用规则组的三种子支持数来扩充样本量。

式\eqref{s0306:joint_criterion}不是$95\%$最优性保证。对于$L_c>0$且$U_c$有限有效的情况，由式\eqref{s0306:bounds}可得
\begin{equation}
 \frac{L_c}{U_c}\leq\frac{L_c}{z_c^*}\leq1,
 \qquad
 \frac{h_{c,a}}{U_c}\leq\frac{h_{c,a}}{z_c^*}
 \leq\min\left\{1,\frac{h_{c,a}}{L_c}\right\}.
 \label{s0306:quality_certificate}
\end{equation}
因此，$L_c/U_c$和$h_{c,a}/U_c$分别提供对应可行收益相对于最优值的下界，而$h_{c,a}/L_c$只是相对于已得参照解的比值。即使$h_{c,a}\geq0.95L_c$，能够据此推出的证书也只是
\begin{equation}
 \frac{h_{c,a}}{z_c^*}\geq
 0.95\frac{L_c}{U_c},
 \label{s0306:threshold_certificate}
\end{equation}
不能把右端擅自替换为$0.95$。缺少有限有效上界时，不报告这类最优性下界。

本节采用预定阈值和描述性统计，未实施正式显著性检验。三个结构种子以及一个公开拓扑的相关干预可以揭示特定条件下的质量与成本变化，但不足以形成独立样本总体上的显著性结论。全部条件仍采用单位容量，其二元可行域可由共享资源的两两冲突表示；较大任务阶数描述资源关联规模，不能独自证明不可约高阶依赖或求解性能优势。
